\section{代理指标与质量门控}
\subsection{关键指标}
除了 SWE-bench 的成功率之外，还需要关注其他指标：人类干预率、平均等待时间、测试覆盖与差异率。这些指标帮助团队判断 Agent 是否真的减轻了工程负担而非增加额外开销。

\begin{importantbox}{指标拆解}
Graham 建议关注四个维度：1) 成功率（pass/fail），2) 人类干预率（多少次需要暂停），3) 运行时间（是否在可接受范围），4) 可解释性（是否输出合理解释）。
\end{importantbox}

\subsection{质量门控节点}
在Agent pipeline 中，质量门控可以设在 Plan、Act 或 Verify 节点：Plan 需要人工 check 生成的步骤，Act 阶段需审查要执行的命令，Verify 阶段则是测试与静态分析。只有所有门控通过后，才允许推送 patch。

\begin{table}[H]
\centering
\begin{tabular}{p{0.3\textwidth} p{0.33\textwidth} p{0.33\textwidth}}
\toprule
门控节点 & 判据 & 响应 \\
\midrule
Plan & 计划的步骤是否合法、有无权限操作 & 人类审查或追加提示 \\
\midrule
Act & 命令是否涉及敏感资源（数据库、云服务） & 暂停并请求审批 \\
\midrule
Verify & 测试/静态分析是否全部通过 & 回滚、增加测试或人工复审 \\
\bottomrule
\end{tabular}
\caption{质量门控在 Agent 流程中的典型布置。}
\end{table}

\begin{warningbox}{不要只看成功率}
高成功率可能意味着 Agent 被限制在低风险任务，实际生产环境还需要综合人类评估与运行成本。
\end{warningbox}

\subsection{本章小结}
\begin{itemize}
    \item 成功率只是其中一个指标，需要结合干预率、运行时间等综合评估。
    \item 质量门控帮助在关键节点介入，防止 Agent 直接发布并引入风险。
\end{itemize}

